\chapter{2025年全国II卷}

\section{单选题}

\begin{problem}
样本数据 \(2\),\(8\),\(14\),\(16\),\(20\) 的平均数为(\quad)

\choices
    {\(8\)}
    {\(9\)}
    {\(12\)}
    {\(18\)}
\end{problem}

\begin{problem}
已知 \(z=1+\i\),则 \(\frac{1}{z-1}=\)(\quad)

\choices
    {\(-\i\)}
    {\(\i\)}
    {\(-1\)}
    {\(1\)}
\end{problem}

\begin{problem}
已知集合 \(A=\{-4,0,1,2,8\}\),\(B=\{x \mid x^3=x\}\),则 \(A\cap B=\)(\quad)

\choices
    {\(\{0,1,2\}\)}
    {\(\{1,2,8\}\)}
    {\(\{2,8\}\)}
    {\(\{0,1\}\)}
\end{problem}

\begin{problem}
不等式 \(\frac{x-4}{x-1}\ge 2\) 的解集是(\quad)

\choices
    {\(\{x \mid -2\le x\le 1\}\)}
    {\(\{x \mid x\le -2\}\)}
    {\(\{x \mid -2\le x<1\}\)}
    {\(\{x \mid x>1\}\)}
\end{problem}

\begin{problem}
在 \(\triangle ABC\) 中,\(BC=2\),\(AC=1+\sqrt3\),\(AB=\sqrt6\),则 \(A=\)(\quad)

\choices
    {\(45^\circ\)}
    {\(60^\circ\)}
    {\(120^\circ\)}
    {\(135^\circ\)}
\end{problem}

\begin{problem}
设抛物线 \(C:y^2=2px\)(\(p>0\))的焦点为 \(F\),点 \(A\) 在 \(C\) 上,过 \(A\) 作 \(C\) 的准线的垂线,垂足为 \(B\). 若直线 \(BF\) 的方程为 \(y=-2x+2\),则 \(|AF|=\)(\quad)

\choices
    {\(3\)}
    {\(4\)}
    {\(5\)}
    {\(6\)}
\end{problem}

\begin{problem}
记 \(S_n\) 为等差数列 \(\{a_n\}\) 的前 \(n\) 项和,若 \(S_3=6\),\(S_5=-5\),则 \(S_6=\)(\quad)

\choices
    {\(-20\)}
    {\(-15\)}
    {\(-10\)}
    {\(-5\)}
\end{problem}

\begin{problem}
已知 \(0<\alpha<\pi\),\(\cos \frac{\alpha}{2}=\frac{\sqrt5}{5}\),则 \(\sin\left(\alpha-\frac{\pi}{4}\right)=\)(\quad)

\choices
    {\(\frac{\sqrt2}{10}\)}
    {\(\frac{\sqrt2}{5}\)}
    {\(\frac{3\sqrt2}{10}\)}
    {\(\frac{7\sqrt2}{10}\)}
\end{problem}

\section{多选题}

\begin{problem}
记 \(S_n\) 为等比数列 \(\{a_n\}\) 的前 \(n\) 项和,\(q\) 为 \(\{a_n\}\) 的公比,\(q>0\). 若 \(S_3=7\),\(a_3=1\),则(\quad)

\choices
    {\(q=\frac12\)}
    {\(a_5=\frac19\)}
    {\(S_5=8\)}
    {\(a_n+S_n=8\)}
\end{problem}

\begin{problem}
已知 \(f(x)\) 是定义在 \(\R\) 上的奇函数,且当 \(x>0\) 时,\(f(x)=(x^2-3)\e^x+2\),则(\quad)

\choices
    {\(f(0)=0\)}
    {当 \(x<0\) 时,\(f(x)=-(x^2-3)\e^{-x}-2\)}
    {\(f(x)\ge 2\) 当且仅当 \(x\ge \sqrt3\)}
    {\(x=-1\) 是 \(f(x)\) 的极大值点}
\end{problem}

\begin{problem}
双曲线 \(C:\frac{x^2}{a^2}-\frac{y^2}{b^2}=1\)(\(a>0\),\(b>0\))的左,右焦点分别是 \(F_1\),\(F_2\),左,右顶点分别为 \(A_1\),\(A_2\). 以 \(F_1F_2\) 为直径的圆与 \(C\) 的一条渐近线交于 \(M\),\(N\) 两点,且 \(\angle NA_1M=\frac{5\pi}{6}\),则(\quad)

\choices
    {\(\angle A_1MA_2=\frac{\pi}{6}\)}
    {\(|MA_1|=2|MA_2|\)}
    {\(C\) 的离心率为 \(\sqrt{13}\)}
    {当 \(a=\sqrt2\) 时,四边形 \(NA_1MA_2\) 的面积为 \(8\sqrt3\)}
\end{problem}

\section{填空题}

\begin{problem}
已知平面向量 \(\bs{a}=(x,1)\),\(\bs{b}=(x-1,2x)\),若 \(\bs{a}\perp(\bs{a}-\bs{b})\),则 \(|\bs{a}|=\)\fillinblank{}.
\end{problem}

\begin{problem}
若 \(x=2\) 是函数 \(f(x)=(x-1)(x-2)(x-a)\) 的极值点,则 \(f(0)=\)\fillinblank{}.
\end{problem}

\begin{problem}
一个底面半径为 \(4\text{cm}\),高为 \(9\text{cm}\) 的封闭圆柱形容器内有两个半径相等的铁球,则铁球半径的最大值为\fillinblank{}\(\text{cm}\).
\end{problem}

\section{解答题}

\begin{problem}
已知函数 \(f(x)=\cos(2x+\varphi)\)(\(0\le \varphi<\pi\)),且 \(f(0)=\frac12\).
\begin{enumerate}
\item 求 \(\varphi\);
\item 设函数 \(g(x)=f(x)+f\left(x-\frac{\pi}{6}\right)\),求 \(g(x)\) 的值域和单调区间.
\end{enumerate}
\end{problem}

\begin{problem}
已知椭圆 \(C:\frac{x^2}{a^2}+\frac{y^2}{b^2}=1\)(\(a>b>0\))的离心率为 \(\frac{\sqrt2}{2}\),长轴长为 4.
\begin{enumerate}
\item 求 \(C\) 的方程;
\item 过点 \((0,-2)\) 的直线 \(l\) 与 \(C\) 交于 \(A\),\(B\) 两点,\(O\) 为坐标原点. 若 \(\triangle OAB\) 的面积为 \(\sqrt2\),求 \(|AB|\).
\end{enumerate}
\end{problem}

\begin{problem}
如图,在四边形 \(ABCD\) 中,\(AB//CD\),\(\angle DAB=90^\circ\),\(F\) 为 \(CD\) 的中点,点 \(E\) 在 \(AB\) 上,\(EF//AD\),\(AB=3AD\),\(CD=2AD\).

将四边形 \(EFDA\) 沿 \(EF\) 翻折至四边形 \(EFD'A'\),使得面 \(EFD'A'\) 与面 \(EFCB\) 所成的二面角为 \(60^\circ\).
\begin{enumerate}
\item 证明:\(A'B//\) 平面 \(CD'F\);
\item 求面 \(BCD'\) 与面 \(EFD'A'\) 所成的二面角的正弦值.
\end{enumerate}

\bitmapfigure[max width=6.0cm]{img/2025/national_paper_2/q17_fig1.png}
\end{problem}

\begin{problem}
已知函数 \(f(x)=\ln(1+x)-x+\frac12x^2-kx^3\),其中 \(0<k<\frac13\).
\begin{enumerate}
\item 证明:\(f(x)\) 在区间 \((0,+\infty)\) 存在唯一的极值点和唯一的零点;
\item 设 \(x_1\),\(x_2\) 分别为 \(f(x)\) 在区间 \((0,+\infty)\) 的极值点和零点.
    \begin{enumerate}
    \item 设函数 \(g(t)=f(x_1+t)-f(x_1-t)\),证明:\(g(t)\) 在区间 \((0,x_1)\) 单调递减;
    \item 比较 \(2x_1\) 与 \(x_2\) 的大小,并证明你的结论.
    \end{enumerate}
\end{enumerate}
\end{problem}

\begin{problem}
甲,乙两人进行乒乓球练习,每个球胜者得 1 分,负者得 0 分. 设每个球甲胜的概率为 \(p\)(\(\frac12<p<1\)),乙胜的概率为 \(q\),\(p+q=1\),且各球的胜负相互独立. 对正整数 \(k\ge 2\),记 \(p_k\) 为打完 \(k\) 个球后甲比乙至少多得 2 分的概率,\(q_k\) 为打完 \(k\) 个球后乙比甲至少多得 2 分的概率.
\begin{enumerate}
\item 求 \(p_3\),\(p_4\)(用 \(p\) 表示);
\item 若 \(\frac{p_4-p_3}{q_4-q_3}=4\),求 \(p\);
\item 证明:对任意正整数 \(m\),\(p_{2m+1}-q_{2m+1}<p_{2m}-q_{2m}<p_{2m+2}-q_{2m+2}\).
\end{enumerate}
\end{problem}

